%! Author = Profesor
%! Date = 06/06/2024


\section{Ejercicios}

\subsection{Ejercicios}


\begin{exercise}
    Los siguientes datos de velocidades de reacción se obtuvieron a \temperatura[c]{25} para la siguiente reacción:

    \begin{equation*}
        \ch{2 A + B -> C}
    \end{equation*}

    \begin{table}[H]
        \centering
        % --- tabla antigua ---
        % \begin{tblr}{Q[c,m]|Q[c,m]|Q[c,m]|Q[c,m]}
            % \hline
            % \SetRow{blue!5!white} Exp. & {\(v_{\text{formación}}\) \\ \unit{\Molar\per\minute}} & {\ch{[A]_i} \\ \unit{\Molar}} & {\ch{[B]_i} \\ \unit{\Molar}} \\
            % \hline
            % 1 & \num{4.0e-4} & \num{0.1} & \num{0.1} \\
            % \hline
            % 2 & \num{1.2e-3} & \num{0.3} & \num{0.3} \\
            % \hline
            % 3 & \num{4.0e-4} & \num{0.1} & \num{0.3} \\
            % \hline
            % 4 & \num{8.0e-4} & \num{0.2} & \num{0.4} \\
            % \hline
        % \end{tblr}
        % --- tabla nueva ---
        \begin{tabular}{@{}cccc@{}}
            \toprule
            \textbf{Exp.} & \textbf{\makecell{\(v_{\text{formación}}\) \\ \unit{\Molar\per\minute}}} & \textbf{\makecell{\ch{[A]_i} \\ \unit{\Molar}}} & \textbf{\makecell{\ch{[B]_i} \\ \unit{\Molar}}} \\
            \midrule
            1 & \num{4.0e-4} & \num{0.1} & \num{0.1} \\[3pt]
            2 & \num{1.2e-3} & \num{0.3} & \num{0.3} \\[3pt]
            3 & \num{4.0e-4} & \num{0.1} & \num{0.3} \\[3pt]
            4 & \num{8.0e-4} & \num{0.2} & \num{0.4} \\
            \bottomrule
        \end{tabular}
    \end{table}


% \begin{figure}[H]
%     \centering
%     \includegraphics[scale=.5]{velocidadesinicialesanaya.png}
%     \caption{\negrita{Datos experimentales}}
%     \label{fig:velocidadesinicialesanaya.png}
% \end{figure}

    \begin{enumerate}[label=\letraminuscula]
        \item Determina los órdenes parciales y el orden total de reacción.
        \item Calcula el valor de la constante cinética
    \end{enumerate}

    \begin{solu}

%        \begin{figure}[H]
%            \centering
%            \includegraphics[scale=0.6]{images/imagesUD5CineticaQuimica/UD5_Problemas/ejerciciogeneral01.png}
%        \end{figure}
    \end{solu}

\end{exercise}

\begin{exercise}
    Escribe la ecuación de velocidad de una reacción: \ch{3 A + 2 B + C -> Productos}, para la cual

    \begin{table}[H]
        \centering
        % --- tabla antigua ---
        % \begin{tblr}{Q[c,m]|Q[c,m]|Q[c,m]|Q[c,m]}
            % \hline
            % \SetRow{blue!5!white} [A] & [B]              & [C]              & Velocidad inicial \\
            % \hline
            % \SI{0.4}{\Molar}          & \SI{0.1}{\Molar} & \SI{0.2}{\Molar} & \(v\)             \\
            % \hline
            % \SI{0.4}{\Molar}          & \SI{0.2}{\Molar} & \SI{0.4}{\Molar} & \(8\ v\)          \\
            % \hline
            % \SI{0.2}{\Molar}          & \SI{0.2}{\Molar} & \SI{0.2}{\Molar} & \(v\)             \\
            % \hline
            % \SI{0.4}{\Molar}          & \SI{0.1}{\Molar} & \SI{0.4}{\Molar} & \(4\ v\)          \\
            % \hline
        % \end{tblr}
        % --- tabla nueva ---
        \begin{tabular}{@{}cccc@{}}
            \toprule
            \textbf{[A]} & \textbf{[B]} & \textbf{[C]} & \textbf{Velocidad inicial} \\
            \midrule
            \SI{0.4}{\Molar} & \SI{0.1}{\Molar} & \SI{0.2}{\Molar} & \(v\) \\[3pt]
            \SI{0.4}{\Molar} & \SI{0.2}{\Molar} & \SI{0.4}{\Molar} & \(8\ v\) \\[3pt]
            \SI{0.2}{\Molar} & \SI{0.2}{\Molar} & \SI{0.2}{\Molar} & \(v\) \\[3pt]
            \SI{0.4}{\Molar} & \SI{0.1}{\Molar} & \SI{0.4}{\Molar} & \(4\ v\) \\
            \bottomrule
        \end{tabular}
    \end{table}

    \begin{solu}
%        \begin{figure}[H]
%            \centering
%            \includegraphics[scale=0.6]{images/imagesUD5CineticaQuimica/UD5_Problemas/ejerciciogeneral02.png}
%        \end{figure}
    \end{solu}
\end{exercise}

\begin{exercise}
    Una reacción de un reactivo A con un reactivo B muestra los siguientes datos de velocidad cuando se estudia a diferentes concentraciones de A y de B:
    \begin{table}[H]
        \centering
        % --- tabla antigua ---
        % \begin{tblr}{Q[c,m]|Q[c,m]|Q[c,m]}
            % \hline
            % \SetRow{blue!5!white} [A] & [B]              & Velocidad inicial                  \\
            % \hline
            % \SI{0.1}{\Molar}          & \SI{0.1}{\Molar} & \SI{4e-4}{\mole\per\litre\per\s}   \\
            % \hline
            % \SI{0.2}{\Molar}          & \SI{0.1}{\Molar} & \SI{1.6e-3}{\mole\per\litre\per\s} \\
            % \hline
            % \SI{0.5}{\Molar}          & \SI{0.1}{\Molar} & \SI{1e-2}{\mole\per\litre\per\s}   \\
            % \hline
            % \SI{0.5}{\Molar}          & \SI{0.5}{\Molar} & \SI{1e-2}{\mole\per\litre\per\s}   \\
            % \hline
        % \end{tblr}
        % --- tabla nueva ---
        \begin{tabular}{@{}ccc@{}}
            \toprule
            \textbf{[A]} & \textbf{[B]} & \textbf{Velocidad inicial} \\
            \midrule
            \SI{0.1}{\Molar} & \SI{0.1}{\Molar} & \SI{4e-4}{\mole\per\litre\per\s} \\[3pt]
            \SI{0.2}{\Molar} & \SI{0.1}{\Molar} & \SI{1.6e-3}{\mole\per\litre\per\s} \\[3pt]
            \SI{0.5}{\Molar} & \SI{0.1}{\Molar} & \SI{1e-2}{\mole\per\litre\per\s} \\[3pt]
            \SI{0.5}{\Molar} & \SI{0.5}{\Molar} & \SI{1e-2}{\mole\per\litre\per\s} \\
            \bottomrule
        \end{tabular}
    \end{table}

    Halla: a) El orden respecto a A, respecto a B y el orden global de la reacción.
    b) El valor de la constante de velocidad.
    c) Su ecuación de velocidad.
    d) Velocidad cuando [A] y [B] valen \SI{0.3}{\Molar}

    \begin{solu}

%        \begin{figure}[H]
%            \centering
%            \includegraphics[scale=0.6]{images/imagesUD5CineticaQuimica/UD5_Problemas/ejerciciogeneral03}
%        \end{figure}

    \end{solu}
\end{exercise}

\begin{exercise}
    A partir de los datos siguientes, calcula la constante de velocidad y comprueba que la reacción es de primer orden:

    \begin{equation*}
        \ch{2 N2O5 -> 4 NO2 + O2}
    \end{equation*}

    \begin{table}[H]
        \centering
        % --- tabla antigua ---
        % \begin{tblr}{Q[c,m]|Q[c,m]|Q[c,m]|Q[c,m]|Q[c,m]|Q[c,m]}
            % \hline
            % \SetColumn{blue!5!white} t (\unit{\minute}) & \num{0}   & \num{67}  & \num{133} & \num{200} & \num{267}  \\
            % \hline
            % \ch{[N2O5]} \SI{e-3}{\Molar}                & \num{2.6} & \num{1.5} & \num{0.9} & \num{0.5} & \num{0.3} \\
            % \hline
        % \end{tblr}
        % --- tabla nueva ---
        \begin{tabular}{@{}cccccc@{}}
            \toprule
            \SetColumn{blue!5!white} t (\unit{\minute}) & \num{0} & \num{67} & \num{133} & \num{200} & \num{267} \\[3pt]
            \ch{[N2O5]} \SI{e-3}{\Molar} & \num{2.6} & \num{1.5} & \num{0.9} & \num{0.5} & \num{0.3} \\
            \bottomrule
        \end{tabular}
    \end{table}

    \begin{solu}

%        \begin{figure}[H]
%            \centering
%            \includegraphics[scale=0.6]{images/imagesUD5CineticaQuimica/UD5_Problemas/ejerciciogeneral04.png}
%        \end{figure}

    \end{solu}

\end{exercise}

\begin{exercise}
    A partir de los datos siguientes, calcula la constante de velocidad y comprueba que la reacción es de segundo orden:

    \begin{equation*}
        \ch{2 HI -> H2 + I2}
    \end{equation*}

    \begin{table}[H]
        \centering
        % --- tabla antigua ---
        % \begin{tblr}{Q[c,m]|Q[c,m]|Q[c,m]|Q[c,m]|Q[c,m]|Q[c,m]}
            % \hline
            % \SetColumn{blue!5!white} t (\unit{\minute}) & \num{0}   & \num{16.67} & \num{33.3} & \num{50} & \num{66.7}  \\
            % \hline
            % \ch{[HI]} \SI{e-3}{\Molar}                  & \num{200} & \num{120}   & \num{61}   & \num{41} & \num{31}   \\
            % \hline
        % \end{tblr}
        % --- tabla nueva ---
        \begin{tabular}{@{}cccccc@{}}
            \toprule
            \SetColumn{blue!5!white} t (\unit{\minute}) & \num{0} & \num{16.67} & \num{33.3} & \num{50} & \num{66.7} \\[3pt]
            \ch{[HI]} \SI{e-3}{\Molar} & \num{200} & \num{120} & \num{61} & \num{41} & \num{31} \\
            \bottomrule
        \end{tabular}
    \end{table}

    \begin{solu}

%        \begin{figure}[H]
%            \centering
%            \includegraphics[scale=0.6]{images/imagesUD5CineticaQuimica/UD5_Problemas/ejerciciogeneral05.png}
%        \end{figure}

    \end{solu}

\end{exercise}

\begin{exercise}
    Una reacción duplica su constante de velocidad cuando pasa de \temperatura[c]{25} a \temperatura[c]{50}.
    Calcula el valor de su energía de activación.

    \begin{solu}

        Aplicamos la ecuación de Arrhenius en su forma integrada entre dos temperaturas:

        \begin{equation*}
            \ln \left( \dfrac{k_2}{k_1} \right) = \dfrac{E_a}{R} \left( \dfrac{1}{T_1} - \dfrac{1}{T_2} \right)
        \end{equation*}

        Con \(T_1 = \temperatura[k]{298}\), \(T_2 = \temperatura[k]{323}\) y \(k_2 = 2\, k_1\) (la constante se duplica), tenemos:

        \begin{equation*}
            \ln(2) = \dfrac{E_a}{R} \left( \dfrac{1}{298} - \dfrac{1}{323} \right)
        \end{equation*}

        Calculamos el paréntesis:

        \begin{equation*}
            \dfrac{1}{298} - \dfrac{1}{323} = 0,0033557 - 0,0030960 = \SI{2.597e-4}{\per\kelvin}
        \end{equation*}

        Despejando \(E_a\):

        \begin{equation*}
            E_a = \dfrac{R \ln(2)}{2,597 \times 10^{-4}} = \dfrac{8,314 \times 0,6931}{2,597 \times 10^{-4}} = \SI{22185}{\joule\per\mole} \approx \energia[kjmol]{22.2}
        \end{equation*}

    \end{solu}
\end{exercise}

\begin{exercise}
    Una reacción tiene una energía de activación de \energia[j]{9e4} y a 30 °C, su constante de velocidad vale \num{8.9e-3}.
    ¿A qué temperatura será \(k = \SI{1e-2}{\litre\per\mole\per\s}\)?

    \begin{solu}

        Aplicamos la ecuación de Arrhenius entre dos temperaturas:

        \begin{equation*}
            \ln \left( \dfrac{k_2}{k_1} \right) = \dfrac{E_a}{R} \left( \dfrac{1}{T_1} - \dfrac{1}{T_2} \right)
        \end{equation*}

        Con \(T_1 = \temperatura[k]{303}\), \(k_1 = \num{8.9e-3}\), \(k_2 = \num{1e-2}\) y \(E_a = \SI{9e4}{\joule\per\mole}\):

        \begin{equation*}
            \ln \left( \dfrac{0,01}{0,0089} \right) = \ln(1,1236) = 0,1165
        \end{equation*}

        \begin{equation*}
            0,1165 = \dfrac{90\,000}{8,314} \left( \dfrac{1}{303} - \dfrac{1}{T_2} \right) = 10\,825,6 \left( \dfrac{1}{303} - \dfrac{1}{T_2} \right)
        \end{equation*}

        Despejando el paréntesis:

        \begin{equation*}
            \dfrac{1}{303} - \dfrac{1}{T_2} = \dfrac{0,1165}{10\,825,6} = 1,076 \times 10^{-5}
        \end{equation*}

        \begin{equation*}
            \dfrac{1}{T_2} = \dfrac{1}{303} - 1,076 \times 10^{-5} = 0,0033003 - 0,0000108 = 0,0032896
        \end{equation*}

        \begin{equation*}
            T_2 = \dfrac{1}{0,0032896} = \SI{304.0}{\kelvin} \approx \temperatura[c]{31}
        \end{equation*}

        La constante de velocidad alcanzará el valor \(k = \SI{1e-2}{\litre\per\mole\per\s}\) a una temperatura de, aproximadamente, \(\SI{304}{\kelvin}\) (\SI{31}{\celsius}).

    \end{solu}
\end{exercise}

\begin{exercise}
    La constante de velocidad de una reacción se expresa según: \(k = A\ e^{\dfrac{-E_a}{R\ T}}\).
    Indica y razona:
    \begin{enumerate}[label=\letraminuscula]
        \item Unidades y significado de cada término.
        \item ¿Qué término se modifica por un catalizador?
        \item La reacción transcurre mediante etapas elementales cuyas energías de activación son: \(E_{a1} = 93\ ; E_{a2} = 148\ ; E_{a3} = 73\).
        ¿Cuál será la etapa más lenta?
        \item ¿Qué
        influencia ejercerá la temperatura?
    \end{enumerate}

    \begin{solu}

        \textbf{a) Unidades y significado de cada término}

        \begin{itemize}
            \item \(k\): constante de velocidad. Su significado es la constante de proporcionalidad entre la velocidad de reacción y las concentraciones de los reactivos elevadas a sus respectivos órdenes parciales. Sus unidades dependen del orden total de la reacción (por ejemplo, \(\si{\s^{-1}}\) para orden 1, \(\si{\litre\per\mole\per\s}\) para orden 2, etc.).
            \item \(A\): factor de frecuencia o factor preexponencial. Está relacionado con la frecuencia de los choques entre las moléculas de reactivos y con la orientación favorable de dichos choques. Tiene las mismas unidades que \(k\).
            \item \(E_a\): energía de activación. Es la energía mínima que deben poseer las moléculas al chocar para que la colisión sea efectiva y se forme el complejo activado. Se expresa en unidades de energía por mol, \(\si{\joule\per\mole}\) o \(\si{\kilo\joule\per\mole}\).
            \item \(R\): constante de los gases ideales, \(\SI{8.314}{\joule\per\mole\per\kelvin}\).
            \item \(T\): temperatura absoluta, en \(\si{\kelvin}\).
        \end{itemize}

        \textbf{b) Término modificado por un catalizador}

        Un catalizador actúa proporcionando un mecanismo de reacción alternativo con una \textbf{menor energía de activación} (\(E_a\)), sin consumirse en el proceso global ni alterar la constante de equilibrio. Por tanto, el término que se modifica es \(E_a\) (disminuye), lo que provoca un aumento del valor de \(k\) y, por tanto, de la velocidad de reacción.

        \textbf{c) Etapa más lenta}

        En un mecanismo de varias etapas, la \textbf{etapa determinante de la velocidad} (la más lenta) es la que presenta la \textbf{mayor energía de activación}, ya que es la que ofrece mayor resistencia a que se produzcan choques efectivos. De los valores dados, \(E_{a2} = 148\) es el mayor, por lo que la \textbf{etapa 2} es la etapa más lenta y, por tanto, la que determina la velocidad global de la reacción.

        \textbf{d) Influencia de la temperatura}

        Según la ecuación de Arrhenius, al aumentar la temperatura \(T\), el exponente \(-E_a/(RT)\) se hace menos negativo, por lo que \(k\) aumenta de forma exponencial. Esto se debe a que, a mayor temperatura, aumenta la energía cinética media de las moléculas y, por tanto, la fracción de choques con energía igual o superior a \(E_a\), es decir, de choques efectivos. En consecuencia, un aumento de la temperatura siempre provoca un \textbf{aumento de la velocidad de reacción}.

    \end{solu}
\end{exercise}

\subsection{Cuestiones}

\begin{question}
    Define el concepto de velocidad de reacción.
    ¿Cuáles son las unidades de la velocidad de reacción? ¿De qué factores depende?

    \begin{solu}
        La velocidad de reacción es la rapidez con la que se produce el consumo de un reactivo o la formación de un producto.
        Definiéndola como la derivada de la concentración respecto del tiempo.

        La unidad de la velocidad de reacción es \si{\mole\per\litre\per\s}

        La velocidad dependerá de: temperatura, concentración, naturaleza de los reactivos, estado físico de las sustancias, presencia o no de catalizadores.
    \end{solu}
\end{question}

\begin{question}
    Define: a) Velocidad de reacción.
    b) Orden de reacción.
    c) Degeneración.
    d) Constante de velocidad.


    \begin{solu}

        \begin{enumerate}[label=\letraminuscula]
            \item La velocidad de reacción es la rapidez con la que un reactivo desaparece o un producto aparece.
            \item Los órdenes parciales de reacción son los exponentes a los que están elevadas las concertaciones en la ecuación de velocidad.
            El orden global de reacción es la suma de los órdenes parciales de reacción.
            \item \negrita{Degeneración}: Es cuando tenemos un orden de reacción cero o algún orden parcial de reacción es cero.
            Cuando tenemos esta situación decimos que la reacción está degenerada respecto de ese reactivo y que su concentración no influirá en la velocidad de reacción.
        \end{enumerate}
    \end{solu}
\end{question}

\begin{question}
    ¿En qué tipo de reacciones es más fácil variar la velocidad variando la concentración de reactivos?

    \begin{solu}

        La velocidad de reacción dependerá de la concentración de los reactivos y por tanto, será más fácil variar la velocidad cuanto más dependa de la concentración, es decir cuando el orden de reacción sea mayor. Por ejemplo, en la ecuación de velocidad de $v = k\ \conc{A}\ \conc{B}^4$ será más fácil variar la velocidad de la reacción variando la concentración de $B$ que la concentración de $A$.

    \end{solu}

\end{question}

\begin{question}
    Un modo de estimar el orden de reacción es mediante las unidades de la constante de velocidad. Para la ecuación:
    \begin{equation*}
        \formula{A -> Productos}
    \end{equation*}
    escribe las correspondientes ecuaciones de velocidad de orden 1 y orden 2 y deduce, para cada caso, las unidades de la constante de velocidad.

    \begin{solu}
        % --- tabla antigua ---
        % \begin{tblr}{|c|c|c|}
            % \hline
            % & \negrita{Orden 1}  & \negrita{Orden 2}              \\
            % \hline
            % Ec. de Velocidad & $ v = k\ \conc{A}$ & $ v = k\ \conc{A}^2$           \\
            % \hline
            % Unidades de $k$  & \unit{\s^{-1}}     & \unit{\litre\mole^{-1}\s^-1{}} \\
            % \hline
        % \end{tblr}
        % --- tabla nueva ---
        \begin{tabular}{@{}ccc@{}}
            \toprule
             & \negrita{Orden 1} & \negrita{Orden 2} \\[3pt]
            Ec. de Velocidad & $ v = k\ \conc{A}$ & $ v = k\ \conc{A}^2$ \\[3pt]
            Unidades de $k$ & \unit{\s^{-1}} & \unit{\litre\mole^{-1}\s^-1{}} \\
            \bottomrule
        \end{tabular}
    \end{solu}
\end{question}

\begin{question}
    Indica cómo afectaría a la velocidad de una reacción:
    \begin{enumerate}[label=\letraminuscula]
        \item El tiempo.
        \item La reducción del volumen del recipiente en una reacción en fase gaseosa.
        \item La presencia de un catalizador.
    \end{enumerate}

    \begin{solu}

        \begin{enumerate}[label=\letraminuscula]
            \item A medida que transcurre el tiempo, los reactivos se van consumiendo y su concentración disminuye; como la velocidad de reacción depende de la concentración de los reactivos, la \textbf{velocidad de reacción disminuye} progresivamente a lo largo del tiempo.
            \item Al reducir el volumen del recipiente en una reacción en fase gaseosa, aumenta la concentración (y la presión parcial) de todas las especies gaseosas, ya que el mismo número de moles ocupa un volumen menor. Al aumentar las concentraciones de los reactivos, \textbf{aumenta la velocidad de reacción}.
            \item Un catalizador proporciona un mecanismo de reacción alternativo con menor energía de activación, facilitando que un mayor número de choques sean efectivos. Por tanto, la presencia de un catalizador \textbf{aumenta la velocidad de reacción}, sin consumirse ni alterar la posición del equilibrio.
        \end{enumerate}

    \end{solu}
\end{question}

\subsection{Ejercicios PAU}

\begin{exerciseEVAU}
    El cloro y el monóxido de nitrógeno reaccionan según la ecuación:
    \begin{centeredequation}
        \ch{2 NO \gas{} + Cl2 \gas{} -> 2 NOCl \gas{}}
    \end{centeredequation}
    Variando las concentraciones iniciales de ambos reactivos se han obtenido los siguientes datos:
    \begin{table}[H]
            \centering
            \caption{}
            \vspace{-5pt}
            \begin{tabular}{@{}lccc@{}}
                \toprule
                \textbf{Expreimento} & \textbf{\( [\ch{NO}]_0 \) (\unit{\M})}       & \textbf{\( [\ch{Cl2}]_0 \) (\unit{\M})}        & \textbf{\( v_0 \) (\unit{\M\per\s})}              \\
                \midrule
                   1      &   0.250      &  0.250  & \num{1.43e-6} \\[3pt]
                2 & 0.5 &  0.250 &  \num{5.72e-6}\\[3pt]
                3 & 0.5 & 0.5 & \num{1.14e-5}\\
                \bottomrule
            \end{tabular}
            \label{tab:}
        \end{table}
    \begin{enumerate}[label=\letraminuscula]
        \item Calcule el orden total de la reacción y escriba la ecuación de velocidad.
        \item Halle la constante de velocidad y determine sus unidades.
        \item Explique cómo afectaría a la velocidad de la reacción, la adición de un catalizador. ¿La
variación de la concentración de qué reactivo afectaría más a la velocidad de reacción? Razónelo.
    \end{enumerate}
\begin{solu}

    \textbf{a) Orden total y ecuación de velocidad}

    Comparamos los experimentos 1 y 2, donde \([\ch{Cl2}]_0\) se mantiene constante y \([\ch{NO}]_0\) se duplica (de 0,250 a 0,5 M):

    \begin{equation*}
        \dfrac{v_2}{v_1} = \dfrac{5,72 \times 10^{-6}}{1,43 \times 10^{-6}} = 4,0 = \left( \dfrac{0,5}{0,250} \right)^n = 2^n \ \Rightarrow \ n = 2
    \end{equation*}

    El orden parcial respecto a \ch{NO} es 2.

    Comparamos ahora los experimentos 2 y 3, donde \([\ch{NO}]_0\) se mantiene constante y \([\ch{Cl2}]_0\) se duplica (de 0,250 a 0,5 M):

    \begin{equation*}
        \dfrac{v_3}{v_2} = \dfrac{1,14 \times 10^{-5}}{5,72 \times 10^{-6}} = 1,99 \approx 2,0 = \left( \dfrac{0,5}{0,250} \right)^m = 2^m \ \Rightarrow \ m = 1
    \end{equation*}

    El orden parcial respecto a \ch{Cl2} es 1. El orden total de la reacción es \(n + m = 2 + 1 = 3\), y la ecuación de velocidad es:

    \begin{equation*}
        v = k\, [\ch{NO}]^2\, [\ch{Cl2}]
    \end{equation*}

    \textbf{b) Constante de velocidad y unidades}

    Sustituyendo los datos del experimento 1:

    \begin{equation*}
        k = \dfrac{v_1}{[\ch{NO}]_1^2\, [\ch{Cl2}]_1} = \dfrac{1,43 \times 10^{-6}}{(0,250)^2 (0,250)} = \dfrac{1,43 \times 10^{-6}}{0,015625} = 9,15 \times 10^{-5}
    \end{equation*}

    (Se obtiene un valor equivalente, \(k \approx 9{,}1$–$9{,}2 \times 10^{-5}\), sustituyendo los datos de los experimentos 2 o 3.)

    Como la reacción es de orden total 3, y \(v\) se mide en \(\unit{\Molar\per\s}\) mientras que las concentraciones se miden en \(\unit{\Molar}\), las unidades de \(k\) son:

    \begin{equation*}
        [k] = \dfrac{\unit{\Molar\per\s}}{\unit{\Molar^3}} = \si{\square\litre\per\square\mole\per\s} = \si{\Molar^{-2}\per\s}
    \end{equation*}

    \textbf{c) Efecto de un catalizador y sensibilidad a cada reactivo}

    Un catalizador proporciona un mecanismo de reacción alternativo con menor energía de activación, por lo que \textbf{aumenta la velocidad de la reacción} (tanto la directa como la inversa, si la reacción fuese reversible), sin consumirse ni modificar la posición del equilibrio, ya que afecta por igual a la cinética de ambos sentidos.

    En cuanto a la sensibilidad a la concentración de cada reactivo, como el orden parcial respecto a \ch{NO} es 2 y respecto a \ch{Cl2} es 1, la velocidad depende de forma cuadrática de \([\ch{NO}]\) y de forma lineal de \([\ch{Cl2}]\). Por tanto, una variación en la concentración de \ch{NO} \textbf{afecta más} a la velocidad de reacción que la misma variación relativa en la concentración de \ch{Cl2}.

\end{solu}
\end{exerciseEVAU}

\begin{exerciseEVAU}
    Si la reacción \ch{A + B -> C} es de segundo orden con respecto al reactivo A y de primer orden con respecto al reactivo B:
    \begin{enumerate}[label=\letraminuscula]
        \item Escriba la ecuación de velocidad de la reacción.
        \item ¿Cuáles son las unidades de la velocidad de reacción (v) y de la constante de
        velocidad (k)?
        \item Justifique cómo variarán la velocidad de reacción (v) y la constante de
        velocidad (k) si se duplica la concentración del reactivo A, manteniendo constante la
        concentración de B.
        \item ¿Cómo afectará a la velocidad de reacción (v) una disminución de la T?

    \end{enumerate}

    \begin{solu}

        \textbf{a) Ecuación de velocidad}

        Como la reacción es de segundo orden respecto a A y de primer orden respecto a B, la ecuación de velocidad es:

        \begin{equation*}
            v = k\, [A]^2\, [B]
        \end{equation*}

        \textbf{b) Unidades de \(v\) y de \(k\)}

        La velocidad de reacción se expresa siempre en unidades de concentración partido de tiempo:

        \begin{equation*}
            [v] = \si{\mole\per\litre\per\s} = \si{\Molar\per\s}
        \end{equation*}

        El orden total de la reacción es \(2 + 1 = 3\), por lo que, despejando \(k\) de la ecuación de velocidad, sus unidades son:

        \begin{equation*}
            [k] = \dfrac{[v]}{[A]^2[B]} = \dfrac{\si{\Molar\per\s}}{\si{\Molar^3}} = \si{\square\litre\per\square\mole\per\s}
        \end{equation*}

        \textbf{c) Efecto de duplicar \([A]\), manteniendo \([B]\) constante}

        La velocidad de reacción es proporcional al cuadrado de \([A]\), por lo que, si \([A]\) se duplica:

        \begin{equation*}
            v' = k\, (2[A])^2\, [B] = 4\, k\, [A]^2\, [B] = 4\, v
        \end{equation*}

        La \textbf{velocidad de reacción se cuadruplica}. Sin embargo, la \textbf{constante de velocidad \(k\) no varía}, ya que \(k\) solo depende de la temperatura (y de la presencia de catalizador), no de las concentraciones de los reactivos.

        \textbf{d) Efecto de una disminución de la temperatura}

        Según la ecuación de Arrhenius, \(k = A\, e^{-E_a/RT}\), al disminuir la temperatura \(T\) disminuye el valor de la constante de velocidad \(k\) (el exponente se hace más negativo), y por tanto \textbf{disminuye la velocidad de reacción}, ya que hay menos moléculas con energía suficiente para superar la energía de activación en cada choque efectivo.

    \end{solu}

\end{exerciseEVAU}

\begin{exerciseEVAU}
    \textbf{2024} La ecuación de velocidad de la reacción entre el monóxido de nitrógeno y el
    dioxígeno es: \( v = K\conc{NO}^2\ \conc{O2} \).
    \begin{enumerate}[label=\letraminuscula]
        \item Determine el orden de reacción respecto a cada reactivo y el orden total
        \item Explique qué ocurriría a la velocidad de la reacción si se duplica la concentración de NO
        y la concentración de \ch{O2} se reduce a la mitad.
    \end{enumerate}

    \begin{solu}

        \textbf{a) Órdenes de reacción}

        En la ecuación de velocidad \(v = k\, [\ch{NO}]^2\, [\ch{O2}]\), el exponente al que está elevada la concentración de cada reactivo es su orden parcial:

        \begin{itemize}
            \item Orden respecto a \ch{NO}: 2
            \item Orden respecto a \ch{O2}: 1
        \end{itemize}

        El orden total de la reacción es la suma de los órdenes parciales: \(2 + 1 = 3\).

        \textbf{b) Efecto de duplicar \([\ch{NO}]\) y reducir a la mitad \([\ch{O2}]\)}

        Sustituyendo los nuevos valores de concentración en la ecuación de velocidad:

        \begin{equation*}
            v' = k\, (2\, [\ch{NO}])^2\, \left( \dfrac{1}{2}\, [\ch{O2}] \right) = k\, \cdot 4\, [\ch{NO}]^2 \cdot \dfrac{1}{2}\, [\ch{O2}] = 2\, k\, [\ch{NO}]^2\, [\ch{O2}] = 2\, v
        \end{equation*}

        Por tanto, la \textbf{velocidad de la reacción se duplicaría}: el efecto de cuadruplicar la velocidad (por duplicar \([\ch{NO}]\), que tiene orden 2) se combina con el efecto de reducirla a la mitad (por reducir a la mitad \([\ch{O2}]\), que tiene orden 1), dando un aumento neto de \(4 \times 0,5 = 2\) veces.

    \end{solu}
\end{exerciseEVAU}

\begin{exerciseEVAU}
    \textbf{2023} Para la reacción:
    \begin{reaction*}
        2 HNO2 + 2 HI -> I2 + 2 NO + 2 H2O
    \end{reaction*}
    se ha calculado experimentalmente la expresión de la velocidad de reacción, la cual es \( v = k [\ch{HNO2}][HI]^2 \).
    Justifique cómo se modifica la velocidad de reacción si:
    \begin{enumerate}[label=\letraminuscula]
    \item Duplicamos la concentración de \ch{HI} y se mantiene constante la de \ch{HNO2}.
    \item Se añade un catalizador positivo al medio.
    \end{enumerate}

    \begin{solu}

        \textbf{a) Efecto de duplicar \([\ch{HI}]\)}

        La ecuación de velocidad es \(v = k\, [\ch{HNO2}]\, [\ch{HI}]^2\), de modo que la velocidad depende del cuadrado de la concentración de \ch{HI}. Si duplicamos \([\ch{HI}]\) manteniendo constante \([\ch{HNO2}]\):

        \begin{equation*}
            v' = k\, [\ch{HNO2}]\, (2\, [\ch{HI}])^2 = 4\, k\, [\ch{HNO2}]\, [\ch{HI}]^2 = 4\, v
        \end{equation*}

        La \textbf{velocidad de reacción se cuadruplica}.

        \textbf{b) Efecto de un catalizador positivo}

        Un catalizador positivo proporciona un mecanismo de reacción alternativo con una energía de activación menor, lo que aumenta la fracción de moléculas capaces de superar la barrera energética en cada choque. Por tanto, \textbf{la velocidad de reacción aumenta}. El catalizador no se consume durante el proceso ni modifica la constante de equilibrio ni las entalpías de reacción; únicamente acelera la velocidad con la que se alcanza el equilibrio.

    \end{solu}
\end{exerciseEVAU}

\begin{exerciseEVAU}
    Para la reacción química en fase gaseosa:
    \begin{equation*}
        \ch{4 HBr + O2 -> 2 H2O + 2 Br2}
    \end{equation*}

    Se sabe que la velocidad de reacción viene dada por la expresión $v = k\ [\ch{HBr}]\ [\formula{O2}]$.
    \begin{enumerate}[label=\letraminuscula]
        \item Indique el orden de reacción.
        \item Para la constante de velocidad, indique sus unidades, sabiendo que las unidades de la velocidad de reacción son \unit{\mole\per\litre\per\s} y las unidades de concentración son \si{\mole\ \litre^{-1}}
    \end{enumerate}

    \begin{solu}

        \begin{enumerate}[label=\letraminuscula]
            \item El orden de reacción sería $1+1 = 2$
            \item Al ser de orden 2 la reacción, las unidades de la constante $k$ serán:

            \si{\litre\per\mole\per\s}
        \end{enumerate}
    \end{solu}
\end{exerciseEVAU}

\begin{exerciseEVAU}
    La ecuación de velocidad de la reacción entre un compuesto A y otro B puede expresarse mediante: $v = K [A][B]^2$.
    Indique cómo debe modificarse la concentración de A para que se mantenga la velocidad de la reacción si la concentración de B se reduce a la mitad.
    \begin{solu}

        Supongamos que se trata de una reacción elemental, y también sabemos que la concentración es directamente proporcional a la velocidad.
        Según la ecuación de velocidad, para que la velocidad de reacción se mantenga con el mismo valor, si la concentración de $B$ se reduce a la mitad, entonces \conc{A} debe ser 4 veces la concentración inicial de $A$.
        Matemáticamente podemos establecer que:
        \begin{equation*}
            v_i = \conc{A}\ \conc{B}^2 \qquad v_f = n\cdot \conc{A}\ (0.5^2\cdot \conc{B}^2)
        \end{equation*}
        Como la condición de es que $v_i = v_f$, igualamos ambas expresiones:
        \begin{gather*}
            \cancelred{\conc{A}}\ \cancelred{\conc{B}^2} = n\cdot \cancelred{\conc{A}}\ (0.5^2\cdot \cancelred{\conc{B}^2}) \separ
            n = \dfrac{1}{0.5^2}\separ
            n = 4
        \end{gather*}

    \end{solu}
\end{exerciseEVAU}

\begin{exerciseEVAU}
    La ecuación de velocidad de la reacción entre el monóxido de nitrógeno y el dihidrógeno es $v = k\;[\ch{NO}]^2\;[\ch{H2}]$ Indique cómo variará la velocidad de la reacción si se duplica la concentración de monóxido de nitrógeno.
    \begin{solu}

        Sabemos que $v = k[\ch{NO}]^2\ [\ch{H2}]$ y, por lo tanto, la velocidad de reacción es directamente proporcional al cuadrado de la concentración de monóxido de nitrógeno, así si se duplica la concentración de \ch{NO} entonces la velocidad aumentará cuatro veces.

        Matemáticamente, también lo podemos establecer de la siguiente manera:
        \begin{equation*}
            v_i = k\;[\ch{NO}]^2\;[\ch{H2}]
        \end{equation*}
        Al duplicar la concentración de \ch{NO} tendremos que:
        \begin{equation*}
            n\ v_i = k\ (2^2\ [\ch{NO}]^2)\ [\ch{H2}]
        \end{equation*}
        Despejamos $n$ que será el que nos dará la información de cuantas veces aumentará la velocidad:
        \begin{equation*}
            v_i = \dfrac{k\ (2^2\ [\ch{NO}]^2)\ [\ch{H2}]}{n}
        \end{equation*}
        Si igualamos ambas expresiones tendremos que:
        \begin{gather*}
            k\;[\ch{NO}]^2\;[\ch{H2}] = \dfrac{k\ (2^2\ [\ch{NO}]^2)\ [\ch{H2}]}{n}
        \end{gather*}
        Despejamos $n$.
        \begin{gather*}
            n = \dfrac{\cancelred{k}\ (2^2\ \cancelred{[\ch{NO}]^2})\ \cancelred{[\ch{H2}]}}{\cancelred{k}\;\cancelred{[\ch{NO}]^2}\;\cancelred{[\ch{H2}]}}\separ
            n = 4
        \end{gather*}

    \end{solu}
\end{exerciseEVAU}

\begin{exerciseEVAU}
    Se encuentra experimentalmente que la ecuación de velocidad para la reacción gaseosa \ch{A + B -> productos}, tiene la forma $v = k\; [A]^2\; [B]$
    \begin{enumerate}[label=\letraminuscula]
        \item Indica el orden de reacción respecto a cada reactivo y el orden total de reacción.
        \item ¿Cuáles son las unidades de la constante de velocidad?
        \item ¿Cómo afecta a la velocidad de la reacción duplicar la concentración de A al tiempo que se reduce a la mitad la de B?
    \end{enumerate}
    \begin{solu}
        \begin{enumerate}[label=\letraminuscula]
            \item El orden de reacción parcial respecto a A es $2$ y de B es $1$ siendo, por tanto, el orden total de la reacción $3$.
            \item Al ser la reacción de orden $3$.
            Las unidades serán:
            \begin{equation}
                \si{\square\litre\per\square\mole\per\s}
            \end{equation}
            \item  Primero duplicar la concentración de A, la velocidad aumentará 4 veces.
            Luego al reducir la concentración de B a la mitad, la reacción disminuirá a la mitad.
            Juntando ambos efectos, tendremos, por tanto, que la reacción aumentará el doble.
        \end{enumerate}
    \end{solu}
\end{exerciseEVAU}

\begin{exerciseEVAU}
     Justifica la certeza o falsedad de las siguientes afirmaciones acerca de la velocidad de una reacción:
    \begin{enumerate}[label=\letraminuscula]
        \item Se modifica cuando se adiciona un catalizador.
        \item Su valor numérico es constante durante todo el tiempo que dura la reacción.
        \item Su valor numérico aumenta al hacerlo la temperatura a la que se realiza la reacción.
        \item Sus unidades pueden ser \si{\mole\s\litre^{-1}}
    \end{enumerate}
    \begin{solu}

        \begin{enumerate}[label=\letraminuscula]
            \item Verdadero, los catalizadores modifican el mecanismo de reacción para diminuir la energía de activación, aumentando así la velocidad de la misma.
            \item Falso, a medida de que el transcurso de la velocidad aumenta, la velocidad disminuye debido a que existen menor reactivos (la velocidad es directamente proporcional a la concentración de los reactivos)
            \item Verdadero, al aumentar la temperatura, las moléculas adquieren mayor velocidad y por lo tanto los choques serán más violentos, alcanzándose más fácilmente la energía de activación aumentando así la velocidad de reacción.
            \item Falso, las unidades de la velocidad siempre son: \si{\mole\per\litre\s}
        \end{enumerate}

    \end{solu}

\end{exerciseEVAU}

\begin{exerciseEVAU}
    La ecuación de velocidad de cierto proceso químico es $v = k\; [A]\; [B]^2$.
    \begin{enumerate}[label=\letraminuscula]
        \item Indica el orden de reacción respecto al reactivo A, respecto al reactivo B y el orden de reacción total.
        \item ¿Cómo variará la velocidad de la reacción si se reduce a la mitad la concentración del reactivo B, manteniendo inalterados el resto de factores?
    \end{enumerate}

    \begin{solu}
        \begin{enumerate}[label=\letraminuscula]
            \item El orden de reacción parcial para A es de $1$ mientras que para B es $2$, siendo el orden total de reacción de $3$.
            \item Si reducimos la concentración de B a la mitad, la velocidad disminuirá 4 veces.
        \end{enumerate}

    \end{solu}
\end{exerciseEVAU}

\newpage

% \subsection{Madrid}

% \myboxcapitulo{} La reacción:
% \begin{center}
%     \ch{CHCl3 (g) + Cl2 (g) -> CCl4 (g) + HCl (g)}
% \end{center}
% es de primer orden con respecto a \ch{CHCl3} y de orden $1/2$ con respecto a \ch{Cl2}. a) Escriba la ecuación de velocidad y determine el orden total de la reacción. b) Deduzca las unidades de la constante de velocidad. c) Justifique cómo afecta a la velocidad de reacción un aumento de volumen a temperatura constante. d) Justifique cómo afecta a la velocidad de reacción un aumento de la temperatura.
% \begin{solu}
%

%
% \end{solu}

% \vspace{6pt}
% \myboxcapitulo{} Para la reacción:
% \begin{center}
%     \ch{2 NO (g) + 2 H2 (g) -> N2 (g) + 2 H2O (g)}
% \end{center}
% el orden parcial de cada reactivo es uno. a) Escriba una expresión para su ecuación de velocidad y calcule el orden total de la reacción. b) Para un valor inicial de \ch{[NO]} y \ch{[H2]} de \concentracion[molar]{0,0025} y \concentracion[molar]{0.075}, respectivamente, la velocidad es $\SI{4,5e-4}{\mole\litre^{-1}\s^{-1}}$.
% \begin{solu}
%

%
% \end{solu}

% \vspace{6pt}
% \myboxcapitulo{} La ecuación de velocidad de la reacción:
% \begin{center}
%     \ch{CO (g) + NO2 (g) -> CO2 (g) + NO (g)}
% \end{center}
% es $v = k\;[NO2]^2$. Justifique si son verdaderas o falsas las siguientes afirmaciones: a) La velocidad de desaparición de ambos reactivos es la misma. b) Las unidades de la constante de velocidad son: \si{\mol\litre\s^-1} c) La velocidad de la reacción aumenta al duplicar la concentración inicial de CO (g). d) En esta reacción en particular, la constante de velocidad no depende de la temperatura, porque la reacción se produce en fase gaseosa.
% \begin{solu}
%

%
% \end{solu}

% \vspace{6pt}
% \myboxcapitulo{} Para la reacción en fase gaseosa:
% \begin{center}
%     \ch{2 NO2 (g) + F2 (g) -> 2 NO2F (g)}
% \end{center}
% la ecuación de velocidad es $v = k\; [\formula{NO2}]\; [\formula{F2}]$. Responda las siguientes cuestiones: a) Indique los órdenes parciales respecto de los reactivos y el orden total de la reacción. b) Razone si es una reacción elemental. c) Determine las unidades de la constante de velocidad. d) Justifique, mediante la ecuación de Arrhenius, cómo afecta un aumento de temperatura a la velocidad de reacción.
% \begin{solu}
%

%
% \end{solu}

% \vspace{6pt}
% \myboxcapitulo{} Se ha llevado a cabo la reacción:
% \begin{center}
%     \ch{A (g)\ + 2 B\ (g) -> 2 C\ (gas)}
% \end{center}
% en dos condiciones experimentales diferentes, obteniéndose la ecuación de velocidad $v = k\; [B]$ y los siguientes valores de energías:
% a) Justifique en cuál de los dos experimentos la reacción es más lenta. b) Explique cómo se modifica la velocidad de la reacción al duplicar la concentración inicial de A. c) Determine el orden total de la reacción y las unidades de la constante de velocidad. d) Justifique cómo afecta a la velocidad de reacción un aumento de temperatura.
% \begin{table}[H]
%     \centering
%     \begin{tblr}{|c|c|c|}
%         \hline
%         \negrita{Experimento} & \negrita{$E_a$ (\si{\kilo\joule\per\mole})} & \negrita{$\Delta H$ (\si{\kilo\joule\per\mole})} \\
%         \hline
%         1                     & \num{2}                                     & \num{-0,3}                                       \\
%         \hline
%         2                     & \num{0,5}                                   & \num{-0,3}                                       \\
%         \hline
%     \end{tblr}
% \end{table}
% \begin{solu}
%

%
% \end{solu}

% \vspace{6pt}
% \myboxcapitulo{} Sabiendo que la ecuación de velocidad $v = k\; [A]^2$ corresponde a la reacción ajustada:
% \begin{center}
%     \formula{ A + 2 B -> C}
% \end{center}
% conteste razonadamente. a) ¿Cuáles son los órdenes parciales de reacción respecto a cada reactivo? ¿Y el orden total de la reacción? b) Deduzca las unidades de la constante de velocidad. c) Indique cómo se modifica la velocidad de la reacción al duplicar la concentración inicial de B. d) Explique cómo afecta a la velocidad de la reacción una disminución de la temperatura.
% \begin{solu}
%

%
% \end{solu}

% \vspace{6pt}
% \myboxcapitulo{} La reacción:
% \begin{center}
%     \ch{3 A (g) + B (g) -> 2 C (g) + D (g)}
% \end{center}
% es de orden 1 respecto de A y de orden 2 respecto de B. a) Escriba la velocidad de la reacción en función de cada especie y justifique si la velocidad de desaparición de B es doble de la velocidad de desaparición de A. b) Obtenga las unidades de la constante de velocidad. c) Razone si la reacción directa es endotérmica sabiendo que la energía de activación es 35 KJ y la de la reacción inversa es \energia[kj]{62}. d) Explique cómo afecta a la velocidad de reacción un aumento de volumen a temperatura constante.
% \begin{solu}
%

%
% \end{solu}

% \vspace{6pt}
% \myboxcapitulo{} Se ha encontrado que la velocidad de la reacción
% \begin{center}
%     \ch{A (g) + 2 B (g) -> C (g)}
% \end{center}
% solo depende de la temperatura y de la concentración de A, de manera que si ésta se triplica, también se triplica la velocidad de la reacción. a) Indica los órdenes de reacción parciales respecto de A y B, así como el orden total. b) Escribe la ley de velocidad. c) Justifica si para el reactivo A cambia más deprisa la concentración que para el reactivo B. d) Explica cómo afecta a la velocidad de reacción una disminución de volumen a temperatura constante.
% \begin{solu}
%

%
% \end{solu}

% \vspace{6pt}
% \myboxcapitulo{} Para la reacción elemental:
% \begin{center}
%     \ch{A(g) + 2 B(g) -> 3 C(g)}
% \end{center}

% a) Escribe la expresión de su ley de velocidad. ¿Cuál es el orden total de la reacción? b) Indica razonadamente cuáles son las unidades de su constante de velocidad. c) ¿Cómo afectará a la velocidad de reacción una disminución de la temperatura a volumen constante? d) Si en un momento determinado se alcanza el estado de equilibrio, indica cómo variarían las cantidades de reactivo si aumentase la presión. ¿Y si se elimina C del medio de reacción?
% \begin{solu}
%

%
% \end{solu}
% \newpage
% \subsection{Comunidad Valenciana}
% \vspace{6pt}
% \myboxcapitulo{} Para el siguiente sistema en equilibrio en fase gaseosa:
% \begin{equation*}
%     \ch{2 NOCl\ (g) <=> 2 NO\ (g) + Cl2\ (g)}
% \end{equation*}
% responda razonadamente a las siguientes cuestiones:
% \begin{enumerate}[label=\letraminuscula]
%     \item Si se extrae del reactor parte del \ch{Cl2\ (g)}, ¿la relación \ch{[NOCl]} / \ch{[NO]} aumenta, disminuye o permanece constante?
%     \item Se observa que al aumentar la temperatura se forma más NOCl. ¿La reacción es exotérmica o endotérmica?
%     \item Si se desea aumentar la cantidad de \ch{NOCl}, manteniendo constante la temperatura, ¿se ha de aumentar o disminuir el volumen del reactor?.
%     \item En un reactor a volumen y temperatura constantes se introducen inicialmente \ch{NOCl} y \ch{Cl2}. Razone si la presión total en el equilibrio será mayor, menor o igual que la inicial.
% \end{enumerate}


% \vspace{6pt}
% \myboxcapitulo{} Considere la reacción:
% \begin{equation*}
%     \ch{3 A\ (g) + 2 B\ (g) -> 2 C\ (g)}
% \end{equation*}
% Se ha observado que, al duplicar la concentración de A, la velocidad de la reacción aumenta cuatro veces mientras que, al disminuir la concentración de B a la mitad, la velocidad disminuye en esa misma proporción. a) Obtenga razonadamente la ley de velocidad de la reacción. b) Cuando las concentraciones iniciales de A y B fueron \SI{0,1}{\molar} y \SI{0,05}{\molar}, respectivamente, la velocidad inicial de la reacción resultó ser \SI{2,82e-4}{\molar\per\s}. Calcule el valor de la constante de velocidad. c) En las condiciones del apartado b), calcule la velocidad de desaparición de A y la velocidad de aparición de C. d) Justifique por qué la velocidad de la reacción aumenta con la temperatura

% \vspace{6pt}
% \myboxcapitulo{} Indique razonadamente si, al disminuir el volumen del reactor a la mitad, manteniendo la temperatura constante, el porcentaje de disociación del \ch{PCl5} aumentará o disminuirá.
% % \negrita{\ul{Datos}}: $R = \R $

% \vspace{6pt}
% \myboxcapitulo{} La cinética de la descomposición del peróxido de hidrógeno, \ch{H2O2}, al reaccionar con el ion yoduro, \formula{I-}, es de primer orden tanto respecto del \formula{H2O2}  como del \formula{I-}. Discuta razonadamente si las siguientes afirmaciones son verdaderas o falsas: a) Un aumento en la concentración de \formula{H2O2} no tiene ningún efecto sobre la velocidad de reacción. b) Al aumentar la temperatura a la que se produce la descomposición del peróxido de hidrógeno, aumenta la velocidad de la reacción. c) La variación en la concentración del ion yoduro afecta más al valor de la velocidad de reacción que la variación de la concentración de \formula{H2O2}. d) La velocidad de la reacción se duplica al duplicar el volumen del reactor, manteniendo constante la temperatura.

% \vspace{6pt}
% \myboxcapitulo{} A una temperatura determinada, se ha estudiado la transformación del NO2 en N2O4 midiendo las velocidades iniciales de la reacción:
% \begin{equation*}
%     \formula{2 NO2\ (g) <=> N2O4\ (g)}
% \end{equation*}
% Se ha determinado que, cuando la concentración inicial de \formula{NO2} es de \SI{0,1}{\molar}, la velocidad inicial de la reacción es \SI{1,45e-4}{\molar\per\s}, mientras que si la concentración inicial de \formula{NO2} es de \SI{0,2}{\molar}, la velocidad inicial de la reacción resulta ser \SI{5,80e-4}{\molar\per\s}. Responda cada una de las siguientes cuestiones:  a) Deduzca la ley de velocidad de la reacción. b) Calcule la constante de velocidad de la reacción en estas condiciones. c) Obtenga la velocidad de desaparición de \formula{NO2} cuando su concentración es \SI{0,14}{\molar} d) Discuta si la velocidad de la reacción aumentará o disminuirá al reducir la temperatura a la cual tiene lugar.

% \vspace{6pt}
% \myboxcapitulo{} En un reactor químico tiene lugar, a \temperatura[c]{800}, la siguiente reacción química:
% \begin{equation*}
%     \formula{CH4\ (g) + H2O\ (g) <=> CO\ (g) + 3 H2\ (g)} \Delta H = \SI{206}{\kilo\joule}
% \end{equation*}

% Responda razonadamente a las siguientes cuestiones: a) Inicialmente, en el recipiente se introducen 1 mol de \formula{CO} y 1 mol de \formula{H2}, manteniendo el volumen y la temperatura constantes. La presión total del recipiente, una vez se alcanza el equilibrio, ¿será mayor, igual o menor que la inicial? Una vez alcanzado el equilibrio: b) Si se quiere que aumente la cantidad de \formula{H2}, ¿habrá que aumentar o disminuir la temperatura? c) Si se quiere que disminuya la cantidad de \formula{CO}, ¿habrá que disminuir o aumentar el volumen? d) Si inyectamos 1 mol de \formula{CO}, manteniendo constantes el volumen y la temperatura, la cantidad de \formula{CH4} aumentará y la cantidad de \formula{H2O} disminuirá. ¿Verdadero o falso?

% \vspace{6pt}
% \myboxcapitulo{} Para la siguiente reacción en fase gaseosa:
% \begin{equation*}
%     \formula{}
% \end{equation*}
% La ecuación de velocidad es $v = k\ [A]^2$ Razone si las siguientes afirmaciones son verdaderas o falsas. a) El reactivo A se consume más deprisa que el reactivo B. b) Las unidades de k son \unit{\litre\; \mole^{-1}\; \min^{-1}}. c) Una vez iniciada la reacción, la velocidad de reacción es constante si la temperatura no varía. d) Al duplicar la concentración de A, a temperatura constante, el valor de la constante de velocidad se cuadruplica.

% \vspace{6pt}
% \myboxcapitulo{} a) La descomposición del pentaóxido de dinitrógeno,
% \begin{equation*}
%     \formula{2 N2O5\ (g) -> 4 NO2\ (g) + O2\ (g)}
% \end{equation*}

% sigue la ecuación de velocidad $v = k\ [\formula{N2O5}]$. Responda las siguientes cuestiones: a1) Compare la velocidad de aparición de \ch{NO2} con la de aparición de \formula{O2}. a2) Indique el orden de reacción total y el orden de reacción respecto del \formula{N2O5}. a3) Indique las unidades de la velocidad de reacción y de la constante de velocidad. a4) Discuta si la constante de velocidad depende de la temperatura a la que se lleva a cabo la reacción. b) Nombre los compuestos siguientes: b1) \formula{K2Cr2O7} b2) \formula{PCl3} b3)\formula{NaClO3}  b4) \formula{Co(OH)2} b5) \formula{FePO4}

% \vspace{6pt}
% \myboxcapitulo{} Discuta razonadamente si las siguientes afirmaciones son verdaderas o falsas.  a) La velocidad para cualquier reacción se expresa en \unit{\mole\; \litre^{-1}\; \s^{-1}} b) Cuando se añade un catalizador a una reacción, esta se hace más exotérmica. c) La velocidad de reacción depende de la temperatura a la que tenga lugar la reacción. d) Para la reacción de segundo orden \formula{A -> B + C}, si la concentración inicial de A es \SI{0.17}{\molar} y la velocidad inicial de la reacción alcanza el valor de \SI{6.8e-3}{\mole\; \litre^{-1}\; \s^{-1}}, la constante de velocidad vale \SI{0.04}{\mole^{-1}\; \litre\; \s^{-1}}.


% \vspace{6pt}
% \myboxcapitulo{} Considere la siguiente reacción:
% \begin{equation*}
%     \formula{CO\ (g) + NO2\ (g) -> CO2\ (g) + NO\ (g)}
% \end{equation*} , cuya ley de velocidad es $v = k\ [\formula{NO2}]^2$. Razone si las siguientes afirmaciones son verdaderas o falsas: a) La velocidad de desaparición del \formula{CO} es igual que la del \formula{NO2}. b) La constante de velocidad no depende de la temperatura porque la reacción se produce en fase gaseosa. c) El orden total de la reacción es cuatro. d) Las unidades de la constante de velocidad serán: \si{\litre\; \mole^{-1}\; \s^{-1}}

% \vspace{6pt}
% \myboxcapitulo{} Considere la reacción:
% \begin{equation*}
%     \formula{2 A + 3 B -> 2 C}
% \end{equation*}
% Se ha observado que al aumentar al doble la concentración de A, la velocidad de la reacción se duplica mientras que al triplicar la concentración de B la velocidad de la reacción aumenta en un factor de 9. Responda razonadamente las siguientes cuestiones:  a) Determine los órdenes de reacción respecto de A y B y escriba la ley de velocidad de la reacción. b) Si en un determinado momento la velocidad de formación de C es \SI{6.12e-4}{\molar\per\s}, calcule la velocidad de la reacción. c) En las mismas condiciones del apartado b), calcule la velocidad de desaparición de B. d) Se ha determinado que cuando las concentraciones iniciales de A y B son \concentracion[m]{0.1} y \concentracion[m]{0.2} respectivamente, la velocidad de la reacción es \SI{2.32e-3}{\molar\per\s}. Calcule la constante de velocidad de la reacción.

% \vspace{6pt}
% \myboxcapitulo{} Considere la reacción:
% \begin{equation*}
%     \formula{A + B -> C}
% \end{equation*}
% Se ha observado que cuando se duplica la concentración de A la velocidad de la reacción se cuadruplica. Por su parte, al disminuir la concentración de B a la mitad, la velocidad de la reacción permanece inalterada. Responda razonadamente las siguientes cuestiones: a) Deduzca el orden de reacción respecto de cada reactivo y escriba la ley de velocidad de la reacción. b) Cuando las concentraciones iniciales de A y B son \concentracion[m]{0.2} y \concentracion[m]{0.1} respectivamente, la velocidad inicial de la reacción alcanza el valor de \SI{3.6e-3}{\molar\; \s^{-1}}. Obtenga el valor de la constante de velocidad. c) ¿Cómo variará la velocidad de la reacción a medida que avance el tiempo? d) ¿Qué efecto tendrá sobre la velocidad de la reacción un aumento de la temperatura a la cual se lleva a cabo?

% \vspace{6pt}
% \myboxcapitulo{} Para la reacción:
% \begin{equation*}
%     \formula{4 NH3\ (g) + 3 O2\ (g) -> 2 N2\ (g) + 6 H2O\ (g)}
% \end{equation*}
% Experimentalmente se determinó que, en un momento dado, la velocidad de formación del \formula{N2} era de \SI{0.27}{\mole\; \litre^{-1}\; \s^{-1}}. Responda a las siguientes cuestiones: a) ¿Cuál era la velocidad de la reacción en ese momento? b) ¿Cuál era la velocidad de formación del agua en ese momento? c) ¿A qué velocidad se estaba consumiendo el \formula{NH3} en ese momento? d) Si la ley de velocidad para esta reacción fuera $v = k\; [\ch{NH3}]^2\; [\ch{O2}]$. ¿Cuáles serían las unidades de la constante de velocidad?

% \newpage

% \subsection{Andalucía}
% \vspace{6pt}
% \myboxcapitulo{} Para la siguiente reacción:
% \begin{center}
%     \formula{2 NO\ (g) + 2 H2\ (g) -> N2\ (g) + 2 H2O\ (g)}
% \end{center}

% la ecuación de velocidad hallada experimentalmente es: $v = k \conc{NO}^2 \conc{H2}$ a) ¿Cuáles son los órdenes parciales de reacción? ¿Y el orden total? b) Si la constante de velocidad para esta reacción a \temperatura[k]{1000} es \SI{6,0e4}{\square\litre\mole^{-2}\s^{-1}}, calcule la velocidad de reacción cuando \conc{NO} = \concentracion[m]{0,015}  y  \conc{H2} = \concentracion[m]{0,035}. c) ¿Cómo afectará a la velocidad de reacción un aumento de la presión, si se mantiene constante la temperatura? Justifique la respuesta.

% \vspace{6pt}
% \myboxcapitulo{} La reacción
% \begin{center}
%     \formula{CO\ (g) + NO2\ (g) -> CO2\ (g) + NO\ (g)}
% \end{center}
% tiene la siguiente ley de velocidad, obtenida experimentalmente: $v = k\ \conc{NO2}^2$ Justifique si son verdaderas o falsas las siguientes afirmaciones: a) La velocidad de desaparición del CO es igual a la velocidad de desaparición del \formula{NO2} b) La constante de velocidad no depende de la temperatura porque la reacción se produce en fase gaseosa. c) El orden total de la reacción es 1 porque la velocidad solo depende de la concentración de \formula{NO2}

% \vspace{6pt}
% \myboxcapitulo{} Indique de forma razonada si las siguientes afirmaciones son verdaderas o falsas: a) La velocidad de una reacción es independiente de la concentración de reactivos. b) La unidad de la constante de velocidad de una reacción de orden uno es \si{\s^{-1}}. c) El uso de catalizadores aumenta la energía de activación de la reacción.

% \vspace{6pt}
% \myboxcapitulo{} La reacción:
% \begin{equation*}
%     \formula{A + 2 B -> C}
% \end{equation*}
% es de orden cero con respecto a A, orden 2 respecto a B y su constante de velocidad vale \SI{0,053}{\litre\mole^{-1}\s{-1}}. a) ¿Cuál es el orden total de la reacción? b) ¿Cuál es la velocidad si las concentraciones iniciales de A y de B son \concentracion[m]{0,48} y \concentracion[m]{0,35}, respectivamente? c) ¿Cómo se modifica la velocidad si la concentración inicial de A se reduce a la mitad?

% \vspace{6pt}
% \myboxcapitulo{} Si la reacción:
% \begin{center}
%     \formula{2 A + B -> C}
% \end{center}
% es de primer orden, justifique si son verdaderas o falsas las siguientes afirmaciones: a) La velocidad de la reacción disminuye al formarse cantidades crecientes de B y C. b) La ecuación de velocidad es $v = k\ \conc{A}^2$. c) Al aumentar la temperatura aumenta la velocidad de la reacción.

% \vspace{6pt}
% \myboxcapitulo{} La reacción elemental:
% \begin{center}
%     \formula{A + B -> C}
% \end{center}
% es de orden 1 para cada reactivo. a) Escriba la ecuación de velocidad correspondiente a dicha reacción. b) A una determinada temperatura la velocidad inicial es de \SI{6,8}{\mole\litre{-1}\s{-1}} y las concentraciones de A y B son \concentracion[molar]{0,17}, calcule la constante de velocidad indicando sus unidades. c) Justifique qué le ocurriría a la velocidad de la reacción si se adiciona un catalizador.

% \vspace{6pt}
% \myboxcapitulo{} La reacción:
% \begin{center}
%     \formula{CO\ (g) + NO2\ (g) ->  CO2\ (g) + NO\ (g)}
% \end{center}
% Tiene la siguiente ecuación de velocidad obtenida experimentalmente: $v = k\ \conc{NO2}^2$. Justifique si son verdaderas o falsas las siguientes afirmaciones: a) La velocidad de desaparición del \ch{CO} es igual a la velocidad de desaparición del \ch{NO2}. b) La constante de velocidad no depende de la temperatura porque la reacción se produce en fase gaseosa. c) El orden total de la reacción es 1 porque la velocidad solo depende de la concentración de \ch{NO2}.

% \vspace{6pt}
% \myboxcapitulo{} Experimentalmente se halla que la reacción:
% \begin{center}
%     \formula{A -> B + C}
% \end{center}
% En fase gaseosa, es de orden 2 respecto de A. a) Escriba la ecuación de velocidad. b) Explique cómo variará la velocidad de reacción si el volumen disminuye a la mitad. c) Calcule la velocidad cuando \conc{A} = \concentracion[m]{0,3}, si la constante de velocidad es $k=\SI{0,36}{\litre\mole^{-1}\s{-1}}$.

% \vspace{6pt}
% \myboxcapitulo{} La reacción:
% \begin{center}
%     \formula{A + 2B + C → D + E}
% \end{center}
% tiene como ecuación de velocidad $v= k\ \conc{A}^2\ \conc{B}$ a) ¿Cuáles son los órdenes parciales de la reacción y el orden total? b) Deduzca las unidades de la constante de velocidad. c) Justifique cuál es el reactivo que se consume más rápidamente

% \vspace{6pt}
% \myboxcapitulo{} Indique verdadero o falso para las siguientes afirmaciones, justificando la respuesta: a) En una reacción del tipo:
% \begin{center}
%     \formula{A + B -> C}
% \end{center}
% el orden total es siempre 2. b) Al aumentar la temperatura a la que se realiza una reacción aumenta siempre la velocidad. c) En un equilibrio la presencia de un catalizador aumenta únicamente la velocidad de la reacción directa.

